\documentclass[letterpaper,oneside]{scrbook}

\usepackage{txfonts}

\usepackage{fontspec}
\usepackage{xunicode}
\usepackage{xltxtra}

\usepackage{fancyvrb}
\usepackage[top=1in,bottom=1in]{geometry}
\usepackage{graphicx}
\usepackage{makeidx}
\usepackage{xcolor}
\usepackage[xetex,breaklinks,colorlinks,linkcolor=black]{hyperref}

\makeindex

\setmainfont[Mapping=tex-text, Numbers=OldStyle]{Palatino LT Std}

\let\primary=\textbf

% Key specs
\def\key#1{{\small$\langle$\addfontfeature{Numbers=Lining}#1\/$\rangle$}}
\let\<=<
\catcode`\<=\active
\def<#1>{\key{#1}}

% Display ‹...› and «...» as text in left and right pointing
% angle brackets. I use «» and ‹› because my terminal doesn't
% display left and right pointing angle brackets properly, and
% Xorg's compose maps don't provide them, anyway.
\catcode`\«=\active
\catcode`\‹=\active
\def‹#1›{$\langle${\itshape#1}$\rangle$}
\def«#1»{$\langle\langle${\itshape#1}$\rangle\rangle$}

% Display |...| as verbatim, teletype text.
\DefineShortVerb{\|}

\makeatletter
\let\:=:
\catcode`\:=\active
\def:{\@ifnextchar:{\coloncoloneq}{\:}}
\def\coloncoloneq#1{\@ifnextchar={$\Coloneqq$\coloncoloneqq}{\:\:}}
\def\coloncoloneqq#1{}

% Create a verbatim {code} environment which highlights strings
% and comments. Several unicode characters are hacked to replace
% the grabbed characters, since we can't escape them in the
% verbatim environment.
\colorlet{comment}{gray}
\colorlet{string}{yellow!80!black!100}
\let\‘=‘
\let\“=“
\catcode`¶=6
\catcode`＃=\active\let＃=\#
\catcode`\#=\active
\catcode`“=\active
\catcode`‘=\active
\def“¶1”{{\color{string}\“¶1”}}%
\def‘¶1’{{\color{string}\‘¶1’}}%
\DefineVerbatimEnvironment{code}{Verbatim}{xleftmargin=2em,gobble=2,%
  codes={\catcode`\#=\active\catcode`\:=\active\catcode`“=\active\catcode`‘=\active},%
  defineactive={%
    \def#{\itshape\color{comment}\let“=\“\let‘=\‘\#}%
    }}
\catcode`\#=6
\catcode`“=12
\catcode`‘=12

% Convenience defs for the various wmii commands, and a few
% others.
\def\wmii{{\tt wmii}}
\def\wiIXmenu{{\tt wi9menu}}
\def\wimenu{{\tt wimenu}}
\def\wmiir{{\tt wmiir}}
\def\ninep{{\addfontfeature{Numbers=Lining}9P}}

\begin{document}
\thispagestyle{empty}
\leavevmode
\vfill

\begin{center}
  \centerline{\includegraphics[width=2in]{../img/wmii.pdf}}

  \vskip 1in

  \LARGE
  The \wmii\ User Guide

  \vskip .5in

  \Large
  Kris Maglione \\[1em]
  \addfontfeature{Numbers=Lining}
  23 May 2009

\end{center}

\vfill

\newpage

\frontmatter

\tableofcontents

\mainmatter

\chapter{Introduction}

\wmii\ is a simple but powerful window manager for the X Window
System. It provides both the classic (“floating”) and tiling
(“managed”) window management paradigms, which is to say, it does
the job of managing your windows, so you don't have to. It also
provides programability by means of a simple file-like
interface, which allows the user to program in virtually any
language he chooses. These basic features have become
indispensible to the many users of \wmii\ and other similar
window managers, but they come at a cost. Though our penchant
for simplicity makes \wmii's learning curve significantly
shorter than most of its competitors, there's still a lot to
learn. The rest of this guide will be devoted to familiarizing
new users with \wmii's novel features and eccentricities, as
well as provide advanced users with an in-depth look at our
customization facilities.

\section{Concepts}

As noted, \wmii\ provides two management styles:

\begin{description}
  \item[Managed] This is the primary style of window management
    in \wmii. Windows managed in this style are automatically
    arranged by \wmii\ into columns. Columns are created and
    destroyed on demand. Individual windows in the column may be
    moved or resized, and are often collapsed or hidden
    entirely. Ad-hoc stacks of collapsed and uncollapsed windows
    allow the user to efficiently manage their tasks. When
    switching from an active to a collapsed window, the active
    window collapses, and the collapsed one effectively takes
    its place.
  \item[Floating] Since some programs aren't designed in ways
    conducive to the managed work flow, \wmii\ also provides the
    classic “floating” window management model. In this model,
    windows float above the managed windows, and may be moved
    freely about. Other than automatic placement of new windows
    and snapping of edges, \wmii\ doesn't manage floating
    windows at all.
  \item[Fullscreen] Fullscreen mode is actually a subset of the
    floating style. Windows may be toggled to and from
    fullscreen mode at will. When fullscreen, windows reside in
    the floating layer, above the managed windows. They have no
    borders or titlebars, and occupy the full area of the
    screen. Other than that, however, they're not special in any
    way. Other floating windows may appear above them, and the
    user can still select, open, and close other windows at
    will.
\end{description}

\subsection{The Filesystem}

All of \wmii's customization is done via a virtual filesystem.
Since the filesystem is implemented in the standardized \ninep\
protocol, it can be accessed in many ways. \wmii\ provides a
simple command-line client, \wmiir, but many alternatives exist,
including libraries for Python, Perl, Ruby, PHP, and C. It can
even be mounted, either by Linux's 9p.ko kernel module or
indirectly via FUSE.

The filesystem that \wmii\ provides is “virtual”, which is to
say that it doesn't reside on disk anywhere. In a sense, it's a
figment of \wmii's imagination. Files, when read, represent
\wmii's current configuration or state. When written, they
perform actions, update the UI, etc. For instance, the directory
|/client/| contains a directory for each window that \wmii\
is currently managing. Each of those directories, in turn,
contains files describing the client's properties (its title,
its views\footnote{Views in \wmii\ are akin to workspaces or
virtual desktops in other window managers, but with some subtle
differences.}, its state). Most files can be written to update
the state they describe. For instance,
|/client/sel/ctl| describes the state of the selected
client. If a client is fullscreen, it contains the line:

\begin{code}
  Fullscreen on
\end{code}

\noindent To change this, you'd update the file with the line
|Fullscreen off| or even |Fullscreen toggle| to toggle
the client's fullscreen state.

The concept of controlling a program via a filesystem derives
from Plan 9, where such interfaces are extensive and well
proven. The metaphor has shown itself to be quite intuitive to
Unix users, once the shock of a “virtual” filesystem wears off.
The flexibility of being able to control \wmii\ from myriad
programming languages, including the standard Unix shell and
even from the command line, is well worth the shock.

\subsection{Views and Tags}

Like most X11 window managers, \wmii\ provides virtual
workspaces. Unlike other window managers, though, \wmii's
workspaces are created and destroyed on demand. Instead of being
sent to a workspace, windows in \wmii\ are tagged with any
number of names. Views are created dynamically from these tags,
and automatically if the user tries to access them.  For
instance, if a window is given the tags ‘foo’ and ‘bar’, the two
views ‘foo’ and ‘bar’ are created, if they don't already exist.
The window is now visible on both of them. Moreover, tags can be
specified as regular expressions. So, a client tagged with {\tt
\verb+/^foo/+} will appear on any view named ‘foo’, ‘foo:bar’,
and so forth. Any time a client is tagged with a matching tag,
or the user opens a matching view, the window is automatically
added to it.

\subsection{The Bar}

\wmii\ provides a general purpose information bar at the top or
bottom of the screen. The bar is divided into a left and a right
section. Each section is made up of buttons, with a single
button spanning the gap between the two sides. Buttons can be
individually styled, and can hold any text content the user
wishes. By convention, the buttons to the left show view names,
and those to the right display status information.

\subsection{The Menus}

\wmii\ includes two simple, external menu programs. The first,
\wimenu, is keyboard-based, and is used to launch programs and
generally prompt the user for input. It provides a list of
completions which are automatically filtered as you type. The
second, \wiIXmenu, is mouse-based, and is generally used to
provide context menus for titlebars and view buttons. Both menus
can be easily launched from shell scripts or the command line,
as well as from more complex scripting languages.

\subsection{The Keyboard}

\wmii\ is a very keyboard friendly window manager. Most actions
can be performed without touching the mouse, including
launching, closing, moving, resizing, and selecting programs.
New keybindings of any complexity can easily be added to handle
any missing functionality, or to simplify any repetitive tasks.

\subsection{The Mouse}

Despite being highly keyboard-accessible, \wmii\ strives to be
highly mouse accessible as well. Windows can be moved or resized
by dragging their window borders. When combined with a key
press, they can be moved, resized, or raised by dragging any
visible portion of the window. Mouse menus are accessed with a
single click and drag. View buttons in the bar and client
titlebars respond to the mouse wheel; view buttons can be
activated by dragging any draggable object (e.g., a file from a
file manager) over them.

\chapter{Getting Started}

This section will walk you through your first \wmii\ startup.
For your first experience, we recommend running \wmii\ in its
own X session, so you can easily switch back to a more
comfortable environment if you get lost. Though you may start
\wmii\ from a session manager in your day to day use, these
instructions will use |xinit|. To begin with, copy this file
to your home directory, so we can open it in your new X session.
Then setup your |~/.xinitrc| as follows:

\begin{code}
  cd

  # Start a PDF viewer with this guide. Use any viewer
  # you're comfortable with.
  xpdf wmii.pdf &

  # Launch wmii
  exec wmii

  # That was easy.
\end{code}

Before you run |xinit|, make sure you know how to switch
between terminals. Depending on your system, your current X
session is probably on terminal 5 or 7. You should be able to
switch between your terminals by pressing
Ctrl-Alt-F$\langle n\rangle$. Assuming that your current X
session is on terminal 7, you should be able to switch between
it and your new session by pressing Ctrl-Alt-F7 and Ctrl-Alt-F8.
Now you should be ready to start \wmii. When you run the
following command, from a terminal, you should be presented with
a new X session, running wmii, with this document open in a PDF
viewer occupying most of the screen. When you're there, come
back to this page and continue. Now, open a terminal and
execute:

\begin{code}
  xinit
\end{code}

\section{Your First Steps}

If everything went according to plan, you should be viewing this
from a nearly empty \wmii\ session. We're going to be using the
keyboard a lot, so let's start with a convention for key
notation. We'll be using the key modifiers Control, Alt, Shift,
and Meta\footnote{The Windows$^{\mbox{\tiny®}}$ key on most
keyboards. The Penguin key, on the more tongue in cheek
varieties.}, which we'll specify as C-, A-, S-, and M-,
respectively. So, <C-S-a> means pressing ‘|a|’ while holding
|Control| and |Shift|.  We'll also express mouse clicks this
way, with <M-Mouse1> signifying a press of the right mouse
button, with the Meta key depressed.  Buttons 4 and 5 are the up
and down scroll wheel directions, respectively.

\subsection{Floating Mode}

Begining with what's familiar to most users, we'll first explore
floating mode. First, we need to select the floating layer.
Press <M-Space>. You should see the titlebar of this window
change color. Now, press <M-Return> to launch a terminal.
The easiest way to drag the terminal around is to press and hold
<M-Mouse1> over the window and simply drag the window
around. You should be able to drag the window anywhere onscreen
without ever releasing the mouse button. As you drag near the
screen edges, you should notice a snap. If you try to drag the
window fully off-screen, you'll find it constrained so that a
portion always remains visible. Now, release the window and move
the mouse toward one of its corners. Press and hold
<M-Mouse3>\footnote{The right button.}. As you drag the
mouse around, you should see the window resized accordingly.

To move the window without the modifier key, move the pointer
over the layout box to the left of its titlebar. You should see
the cursor change. Now, simply click and drag. To resize it,
move the pointer toward the window's edge until you see the
cursor change, and again, click and drag. Now, to close the
window, move the mouse over the windows titlebar, press and hold
<Mouse3>, select |Delete|, and release it. You should
see this window's titlebar return to its original color,
indicating that it's regained focus.

\subsection{Managed Mode}

Now, for the fun part. We'll start exploring managed mode by
looking at the basics of columns. In the default configuration,
columns have three modes:

\begin{description}
  \item[Stack] <M-s> The default mode for new columns. Only one window
    is fully visible per column at once. The others only display
    their title bars. When new windows are added to the column,
    the active window collapses, and the new one takes its
    place. Whenever a collapsed client is selected, the active
    window is collapsed to take its place.
  \item[Max] <M-m> Like stack mode, but the titlebars of collapsed
    clients are hidden.
  \item[Default] <M-d> Multiple uncollapsed windows may be visible at
    once. New windows split the space with the other uncollapsed
    windows in their vicinity. Windows may still be collapsed by
    shrinking them to the size of their titlebars. At this
    point, the behavior of a stack of collapsed and uncollapsed
    clients is similar to that of stack mode.
\end{description}

Before we open any new windows in managed mode, we need to
explore the column modes a bit. Column modes are activated with
the key bindings listed above. This column should be in stack
mode now. Watch the right side of the titlebar as you press
<M-m> to enter max mode. You should see an indicator appear.
This tells you the number of hidden windows directly above and
below the current window, and its position in that stack. Press
<M-d> to enter default mode. Now we're ready to open another
client. Press <M-Return> to launch another terminal. Now,
press <M-S-l> to move the terminal to a new column to the
right of this one. Once it's there, press <M-Return> two
more times to launch two more terminals. Now that you have more
than one window in a column, cycle through the three column
modes again until they seem familiar.

\subsection{Keyboard Navigation}

To begin, switch back to default mode. The basic keyboard
navigation keys, <M-h>, <M-j>, <M-k>, and <M-l>,
derive from vi, and represent moving left, down, up, and right
respectively. Try selecting each of the four windows currently
visible on screen. Notice that navigation wraps from one side of
the screen to the other, and from the top to the bottom. Now,
return to the write column, switch to stack mode, and select
each of the three terminals again. Do the same in max mode,
paying careful attention to the indicator to the right of the
titlebar.

Now that you can select windows, you'll want to move them
around. To move a window, just add the Shift key to the
direction keys. So, to move a window left, instead of <M-h>,
type <M-S-h>. Now, experiment with moving windows, just as
you did with navigating them, in each of the three column modes.
Once you're comfortable with that, move a window to the floating
layer. Since we toggled between the floating and managed layers
with <M-Space>, we'll move windows between them with
<M-S-Space>. Try moving some windows back and forth until it
becomes familiar. Now, move several windows to the floating
layer and try switching between them with the keyboard. You'll
notice that <M-h> and <M-l> don't function in the
floating layer. This is for both historical and logistical
reasons. <M-j> and <M-k> cycle through floating windows
in order of their most recent use.

\subsection{Mouse Navigation}

\wmii\ uses the “sloppy focus” model, which is to say, it focuses
windows when the mouse enters them and when you click them. It
focuses windows only when you select them with the keyboard,
click their titlebars, or press click them with <M-Mouse2>.
Collapsed windows may be opened with the mouse by clicking their
titlebars. Moving and resizing floating windows should be
largely familiar, and has already been covered. The same can't
be said for managed windows.

Let's begin working with the mouse in the managed layer. Return
to a layout with this document in a column on the left, and
three terminals in a column to the right. Switch the right
column to default mode. Now, bring the mouse to the top of the
third terminal's titlebar until you see a resize cursor. Click
and drag the titlebar to the very top of the screen. Now, move
the cursor to the top of the second terminal's titlebar and drag
it to the very bottom of the screen. Press <M-d> to restore the
terminals to their original sizes. Now, click and hold the
layout box of the second terminal. Drag it to the middle of the
terminal's window and release. Click and hold the layout box of
the third terminal and drag it to the middle of the first
terminal's window. Finally, drag the first terminal's layout box
to halfway down this window. <M-Mouse1> works to the same
effect as dragging the layout box, but allows you to click
anywhere in the window.

Now that you've seen the basics of moving and dragging windows,
let's move on to columns. Click and drag the border between the
two columns. If that's a difficult target to click, there's a
triangle at the top of the division between the two columns that
you can click and drag as well. If that's still too hard a
target, try using <M-Mouse3>, which works anywhere and provides
much richer functionality.

\section{Running Programs}

You've already seen the convenient key binding to launch a
terminal, but what about other programs? To get a menu of all of
the executables in your path, type <M-p>. This should replace
the bar at the bottom of the screen with a prompt, followed by a
string of completions. Start typing the name of a program that
you want to open. You can press <Tab> and <S-Tab> to cycle
through the completions, or you can just press <Return> to
select the first one. If you want to execute a more complex
command, just type it out and press <Return>. If you want to
recall that command later, use \wimenu's history. Start typing
the command you want and then press <C-p> until you come to it.

When you're done with a program, you'll probably want an easy
way to close it. The first way is to ask the program to close
itself. Since that can be tedious (and sometimes impossible),
\wmii\ provides other ways. As mentioned, you can right click
the titlebar and select |Delete|. If you're at the keyboard,
you can type <M-S-c>. These two actions cause \wmii\ to ask
nicely that the program exit. In those sticky cases where the
program doesn't respond, \wmii\ will wait 10 seconds before
prompting you to kill the program. If you don't feel like
waiting, you can select |Kill| from the window's titlebar
menu, in which case \wmii\ will forcefully and immediately kill
it. Beware, killing clients is a last resort. In cases where the
same program opens multiple windows, killing one will kill them
all—without warning.

\section{Using Views}

As already noticed, \wmii's concept of virtual workspaces is
somewhat unique, so let's begin exploring it. Open up a terminal
and press <M-S-2>. You should see a new button on the bar at the
bottom of the screen. When you click it, you should see your
original terminal. Press <M-1> to come back here. Now, press
<M-3>, and <M-1> again to return here once more. Notice that the
views were created when needed, and destroyed when no longer
necessary. If you want to select a view with a proper name, use
<M-t> and enter the name. Other than the dynamic creation of
views, this is still similar to the familiar X11 workspace
model. But that's just the begining of \wmii's model. Open a new
terminal, and type:

\begin{code}
  echo ‘Hello world!’
\end{code}

\noindent Now, type <M-S-t>. In the menu that appears, enter
|1+2+3|. Now, visit the views |1|, |2|, and |3|, and you'll see
the client on each. To remove a tag, type <M-S-t> again, and
this time enter |-2|. You'll notice that the client is no longer
on the |2| view. Finally, tag names needn't be discrete,
ordinary strings. They can also be regular expressions. Select
the terminal again, and enter |+/^5/|. Now, switch to the |5|
view.  Now try the |6| view. Finally, type <M-t> and enter |50|
to check the |50| view. Clients tagged with regular expressions
are attached to any matching views when they're created. So,
when you switch to an empty view, or tag a client with a new
tag, any clients with matching regular expressions are
automatically added to it. When all explicitly tagged clients
disappear from the view, and it's no longer visible, clients
held there by regular expressions are automatically removed.

\section{Learning More}

For full tables of the standard key bindings, and descriptions
of the precise semantics of the topics discussed above, you
should refer to \wmii's |man| pages.

\chapter{Customizing \wmii}

There are several configuration schemes available for \wmii. If
you're only looking to add basic key bindings, status monitors,
{\it et cetera}, you should have no trouble modifying the stock
configuration for your language of choice. If you're looking for
deeper knowledge of \wmii's control interface, though, this
section is for you. We'll proceed by building a configuration
script in POSIX |sh| syntax, and move on to a discussion of the
higher level constructs in the stock configuration scripts.

\section{Events}

The \wmii\ control interface is largely event driven. Each event
is represented by a single, plain-text line written to the
|/event| file. You can think of this file as a named pipe. When
reading it, you won't receive an EOF\footnote{End of File} until
\wmii\ exits. Moreover, any lines written to the file will be
transmitted to all of its readers. Notable events include key
presses, the creation and destruction of windows, and changes of
focus and views.

We'll start building our configuration with an event processing
framework:

\begin{code}
  «Event Loop» ::=
  # Broadcast a custom event
  wmiir xwrite /event Start wmiirc

  # Turn off globbing
  set -f
  # Open /event for reading
  wmiir read /event |
  # Read the events line by line
  while read line; do
      # Split the line into words, store in $@
      set -- $line
      event=$1; shift
      line = "$(echo $line | sed ‘s/^[^ ]* //’ | tr -d ‘\n’)"
      # Process the event
      case $event in
      Start) # Quit when a new instance starts
          [ $1 = wmiirc ] && exit;;
      «Event Handlers»
      esac
  done
\end{code}

Now, we need to consider which types of events we'll need to
handle:

\begin{code}
  «Event Handlers» ::=
  «View Button Events»
  «Urgency Events»
  «Unresponsive Clients»
  «Notice Events»
  «Key Events»
  «Client Menu Events»
  «Tag Menu Events»
\end{code}

\section{Bar Items}

The bar is described by the files in the two directories |/lbar/| and
|/rbar/| for buttons on the left and right side of the bar,
respectively. The format of the files is:

\begin{code}
  ‹Color Tuple› ‹Label›
\end{code}

although the color tuple may be elided in cases where the label
doesn't match its format.

A ‹Color Tuple› is defined as:

\begin{code}
  ‹tuple› ::= ‹foreground color› ‹background color› ‹border color›
  ‹color› ::= ＃‹6 character RGB hex color code›
\end{code}

Let's define our basic theme information now:

\begin{code}
  «Theme Definitions» ::=
  normcolors=‘＃000000 ＃c1c48b ＃81654f’
  focuscolors=‘＃000000 ＃81654f ＃000000’
  background=‘＃333333’
  font=‘drift,-*-fixed-*-*-*-*-9-*-*-*-*-*-*-*’
\end{code}

\subsection{View Buttons}

With a basic understanding of bar items in mind, we can write
our view event handlers:

\index{events!CreateTag}
\index{events!DestroyTag}
\index{events!FocusTag}
\index{events!UnfocusTag}
\begin{code}
  «View Button Events» ::=
  CreateTag)  # CreateTag ‹Tag Name›
      echo $normcolors $1 | wmiir create /lbar/$1;;
  DestroyTag) # DestroyTag ‹Tag Name›
      wmiir rm /lbar/$1;;
  FocusTag)   # FocusTag ‹Tag Name›
      wmiir xwrite /lbar/$1 $focuscolors $1;;
  UnfocusTag) # UnfocusTag ‹Tag Name›
      wmiir xwrite /lbar/$* $normcolors $1;;
\end{code}

\subsection{Urgency}

\index{events!UrgentTag|(}
\index{events!NotUrgentTag|(}
Windows can specify that they require attention, and in X11
parlance, this is called urgency. When a window requests
attention as such, or declares that it's been satisfied, \wmii\
broadcasts an event for the client and an event for each view
that it belongs to, and fills in the client's layout box. It's
the job of a script to decide how to handle it above and beyond
that. The standard scripts simply mark urgent views with an
asterisk:

\begin{code}
  «Urgency Events» ::=
  # The urgency events are ‘Client’ events when the program
  # owning the window sets its urgency state. They're ‘Manager’
  # events when wmii or the wmii user sets the state.
  UrgentTag)    # UrgentTag ‹‘Client’ or ‘Manager’› ‹Tag Name›
      wmiir xwrite /lbar/$2 *$2;;
  NotUrgentTag) # NotUrgentTag ‹‘Client’ or ‘Manager’› ‹Tag Name›
      wmiir xwrite /lbar/$2 $2;;
\end{code}
\index{events!UrgentTag|)}
\index{events!NotUrgentTag|)}

\subsection{Notices}

The standard scripts provide a custom Notice event for
displaying status information. The events appear in the long bar
between the left and right sides for five seconds.

\begin{code}
  «Notice Events» ::=
  Notice)
      wmiir xwrite /rbar/!notice $line
      kill $xpid 2>/dev/null # Let's hope this isn't reused...
      { sleep 5; wmiir xwrite /rbar/!notice ‘ ’; } &
      xpid = $!;;
\end{code}

\section{Keys}

Now to the part you've no doubt been waiting for: binding keys.
When binding keys, you need to be aware of two files, |/keys|
and |/event|. The former defines which keys \wmii\ needs to
grab, and the latter broadcasts the events when they're pressed.
Key names are specified as a series of modifiers followed by a
key name, all separated by hyphens. Valid modifier names are
|Control|, |Shift|, |Mod1| (usually Alt), |Mod2|, |Mod3|, |Mod4|
(usually the Windows® key), and |Mod5|. Modifier keys can be
changed via |xmodmap(1|), which is beyond the scope of this
discussion. Key names can be detected by running |xev| from a
terminal, pressing the desired key, and looking at the output
(it's in the parentheses, after the keysym). A \wmii-specific
utility is forthcoming.

Examples of key bindings:

\begin{description}
  \item[Windows® key + Capital A] |Mod4-Shift-A|
  \item[Control + Alt + Space]    |Mod1-Control-Space|
\end{description}

Now, let's bind the keys we plan on using:

\begin{code}
  «Bind Keys» ::=
  {
  cat <<!
  Mod4-space
  Mod4-d
  Mod4-s
  Mod4-m
  Mod4-a
  Mod4-p
  Mod4-t
  Mod4-Return
  Mod4-Shift-space
  Mod4-f
  Mod4-Shift-c
  Mod4-Shift-t
  Mod4-h
  Mod4-j
  Mod4-k
  Mod4-l
  Mod4-Shift-h
  Mod4-Shift-j
  Mod4-Shift-k
  Mod4-Shift-l
  !
  for i in 1 2 3 4 5 6 7 8 9 0; do
      echo Mod4-$i
      echo Mod4-Shift-$i
  done
  } | wmiir write /keys
\end{code}

and lay a framework for processing their events:

\begin{code}
  «Key Events» ::=
  Key) # Key ‹Key Name›
      case $1 in
      «Motion Keys»
      «Client Movement Keys»
      «Column Mode Keys»
      «Client Command Keys»
      «Command Execution Keys»
      «Tag Selection Keys»
      «Tagging Keys»
      esac;;
\end{code}

\section{Click Menus}

Sometimes, you have your hand on the mouse and don't want to
reach for the keyboard. To help cope, \wmii\ provides a
mouse-driven, single-click menu. The default configuration uses
it for client and tag menus.

\begin{code}
  «Click Menu Initialization» ::=
  clickmenu() {
      if res=$(wmii9menu -- “$@”); then eval “$res”; fi
  }
\end{code}

\section{Control Files}

Most filesystem objects, including the root directory, have
control files, named |ctl|. The first line of most control files
is the canonical name of the directory they reside in, which
comes in handy for the special |sel/| directories, which are
aliases for the currently selected object of a group. The
following lines represent properties of the object. Control
files may be written to, in similar syntax to the values that
can be read, to update those properties. For instance, if a file
contains:

\begin{code}
  Fullscreen on
\end{code}

\noindent either of the following, when written to the file,
will disable the |Fullscreen| state:

\begin{code}
  Fullscreen off
  Fullscreen toggle
\end{code}

\section{Clients}

Clients are represented by directories under the |/client/|
tree. Subdirectory names represent the client's X11 window ID.
The special |sel/| directory represents the currently selected
client. The files in these directories are:

\begin{description}
  \item[ctl] The control file. The properties are:
    \begin{description}
      \item[Fullscreen] The client's fullscreen state. When
        |on|, the client is displayed fullscreen on all of its
        views. Possible values are |on|, |off|, and |toggle|.
      \item[Urgent] The client's urgency state. When |on|, the
        client's layout box will be highlighted. Possible values
        are |on|, |off|, and |toggle|.
      \item[kill] When written, the window is closed politely,
        if possible.
      \item[slay] When written, the client is killed peremptorily.
    \end{description}
  \item[props] The client's window class (the X11 |WM_CLASS|
    property) and title string, separated by colons. This file
    is not writable.
  \item[label] The client's window title. May be written to
    change the client's title.
  \item[tags] The client's tags. Tag names are separated by |+|
    signs. Tags beginning and ending with |/| are treated as
    regular expressions. If the written value begins with a |+|
    or a |-|, the tags are updated rather than overwritten. Tag
    names which directly follow a |-| sign are removed rather
    than added. Regular expression tags which directly follow a
    minus sign are treated as exclusion expressions. For
    example, the tag string |+/foo/-/food/| will match the tag
    |foobar|, but not the tag |foodstand|.
\end{description}

\subsection{Key Bindings}

To control clients, we'll add the following key bindings:

\begin{code}
  «Client Command Keys» ::=
  Mod4-Shift-c) wmiir xwrite /client/sel/ctl kill;;
  Mod4-f) wmiir xwrite /client/sel/ctl Fullscreen toggle;;
\end{code}

And to manage their tags, we'll need:

\begin{code}
  «Tagging Keys» ::=
  Mod4-Shift-t)
    # Get the selected client's id
    c=$(wmiir read /client/sel/tag | sed 1q)
    # Prompt the user for new tags
    tags=$(wmiir ls /tag | sed ‘s,/,,; /sel/d’ | wimenu)
    # Write them to the client
    wmiir xwrite /client/$c/tags $tag;;
  Mod4-Shift-[0-9])
    wmiir xwrite /client/sel/tags ${2＃＃*-};;
\end{code}

\subsection{Click Menus}

\index{events!ClientMouseDown}
\begin{code}
  «Client Menu Events» ::=
  ClientMouseDown) # ClientMouseDown ‹Client ID› ‹Button›
    [ $2 = 3 ] && clickmenu \
      “Delete:xwrite /client/$1/ctl kill” \
      “Kill:xwrite /client/$1/ctl slay” \
      “Fullscreen:/client/$1/ctl Fullscreen on”
\end{code}

\subsection{Unresponsive Clients}

\index{events!UnresponsiveClient|(}
When \wmii\ tries to close a window, it waits 8 seconds for the
client to respond, and then lets its scripts decide what to do
with it. The stock scripts prompt the user for input:

\begin{code}
  «Unresponsive Clients» ::=
  UnresponsiveClient) # UnresponsiveClient ‹Client ID›
  {
      # Use wihack to make the xmessage a transient window of
      # the problem client. This will force it to open in the
      # floaing layer of whatever views the client is attached to
      resp=$(wihack -transient $1 \
             xmessage -nearmouse -buttons Kill,Wait -print \
             “The following client is not responding.” \
             “What would you like to do?$(echo)” \
             $(wmiir read /client/$1/label))
      [ $resp = Kill ] && wmiir xwrite /client/$1/ctl slay
  } &;;
\end{code}
\index{events!UnresponsiveClient|)}

\section{Views}

Views are represented by directories under the |/tag/| tree. The
special |sel/| directory represents the currently selected
client. The |sel| tag is treated similarly elsewhere. The files
in these directories are:

\begin{description}
  \item[ctl] The view's control file. The properties are:
  \item[select ‹Area›] Select the column ‹Area›, where
    ‹Area› is a 1-based column index, or |~| for the floating
    area.
  \item[select ‹Area› ‹Client Index›] Select the column ‹Area›, and
    the ‹Client Index›th client.
  \item[select client ‹Client ID›] Select the client with the
    X11 window ID ‹Client ID›.
  \item[select ‹Direction›]
    Select the client in ‹Direction› where ‹Direction› may be
    one of ‹up $\wedge$ down $\wedge$ left $\wedge$ right›.
  \item[send client ‹Client ID› ‹Area›] Send ‹Client ID› to
    ‹Area›. ‹Area› may be |sel| for the selected area, and
    |client ‹Client ID›| may be |sel| for the currently selected
    client.
  \item[send client ‹Client ID› ‹Direction›]
    Send ‹Client ID› to a column or position in its column in
    the given direction.
  \item[send client ‹Client ID› toggle] If ‹Client ID› is
    floating, send it to the managed layer. If it's managed,
    send it to the floating layer.
  \item[swap client ‹Client ID› \ldots] The same as the |send|
    commands, but swap ‹Client ID› with the client at the given
    location.
  \item[colmode ‹Area› ‹Mode›] Set ‹Area›'s mode to ‹Mode›,
    where ‹Mode› is a string of values similar to tag
    specifications. Values which may be added and removed are as
    follows for managed areas:

    \begin{description}
      \item[stack] One and only one client in the area is
        uncollapsed at any given time. When a new client is
        selected, it is uncollapsed and the previously selected
        client is collapsed.
      \item[max] Collapsed clients are hidden from view
        entirely. Uncollapsed clients display an indicator
        {\it‹n›/‹m›}, where ‹m› is the number of collapsed
        clients directly above and below the client, plus one,
        and ‹n› is the client's index in the stack.
    \end{description}

    For the floating area, the values are the same, except that
    in |max| mode, floating clients are hidden when the managed
    layer is selected.
  \item[grow ‹Frame› ‹Direction› {[‹Amount›]}] Grow ‹Frame› in
    the given direction, by ‹Amount›. ‹Amount› may be any
    integer, positive or negative. If suffixed with |px|,
    it specifies an exact pixel amount, otherwise it specifies a
    “reasonable increment”. Defaults to 1.

    ‹Frame› may be one of:
    \begin{itemize}
      \item client ‹Client ID›
      \item ‹Area› ‹Client Index›
    \end{itemize}
  \item[nudge ‹Frame› ‹Direction› {[‹Amount›]}] The same as
    |grow|, but move the client in ‹Direction› instead of
    resizing it.
\end{description}

\subsection{Key Bindings}

We'll use the following key bindings to interact with views:

\begin{code}
  «Motion Keys» ::=
  Mod4-h) wmiir xwrite /tag/sel/ctl select left;;
  Mod4-l) wmiir xwrite /tag/sel/ctl select right;;
  Mod4-k) wmiir xwrite /tag/sel/ctl select up;;
  Mod4-j) wmiir xwrite /tag/sel/ctl select down;;
  Mod4-space) wmiir xwrite /tag/sel/ctl select toggle;;

  «Client Movement Keys» ::=
  Mod4-Shift-h) wmiir xwrite /tag/sel/ctl send sel left;;
  Mod4-Shift-l) wmiir xwrite /tag/sel/ctl send sel right;;
  Mod4-Shift-k) wmiir xwrite /tag/sel/ctl send sel up;;
  Mod4-Shift-j) wmiir xwrite /tag/sel/ctl send sel down;;
  Mod4-Shift-space) wmiir xwrite /tag/sel/ctl send sel toggle;;

  «Column Mode Keys» ::=
  Mod4-d) wmiir xwrite /tag/sel/ctl colmode sel -stack-max;;
  Mod4-s) wmiir xwrite /tag/sel/ctl colmode sel stack-max;;
  Mod4-m) wmiir xwrite /tag/sel/ctl colmode sel stack+max;;
\end{code}

\subsection{Click Menus}

\index{events!LeftBarMouseDown}
\begin{code}
  «Tag Menu Events» ::=
  LeftBarMouseDown) # LeftBarMouseDown ‹Button› ‹Bar Name›
    [ $1 = 3 ] && clickmenu \
      “Delete:delete_view $2”
\end{code}

\section{Command and Program Execution}

Perhaps the most important function we need to provide for is
the execution of programs. Since \wmii\ users tend to use
terminals often, we'll add a direct shortcut to launch one.
Aside from that, we'll add a menu to launch arbitrary programs
(with completions) and a separate menu to launch wmii specific
commands.

We use |wmiir setsid| to launch programs with their own session
IDs to prevent untoward effects when this script dies.

\begin{code}
  «Command Execution Initialization» ::=
  terminal() { wmiir setsid xterm “$@” }
  proglist() {
      IFS=: set -- $1
      find -L $@ -maxdepth 1 -perm /111 | sed ‘s,.*/,,’ | sort | uniq
      unset IFS
  }
\end{code}

\subsection{Key Bindings}
\begin{code}
  «Command Execution Keys» ::=
  Mod4-Return) terminal & ;;
  Mod4-p) eval exec wmiir setsid "$(proglist $PATH | wimenu)" &;;
  Mod4-a) {
      set -- $(proglist $WMII_CONFPATH | wimenu)
      prog = $( (PATH=$WMII_CONFPATH which $1) ); shift
      eval exec $prog “$@”
  } &;;
\end{code}

\section{The Root}

The root filesystem contains the following:

\begin{description}
  \item[ctl] The control file. The properties are:
    \begin{description}
      \item[bar on ‹top $\wedge$ bottom›] Controls where the bar
        is shown.
      \item[bar off] Disables the bar entirely.
      \item[border] The border width, in pixels, of floating
        clients.
      \item[colmode ‹Mode›] The default column mode for newly
        created columns.
      \item[focuscolors ‹Color Tuple›] The colors of focused
        clients.
      \item[normcolors ‹Color Tuple›] The colors of unfocused
        clients and the default color of bar buttons.
      \item[font ‹Font›] The font used throughout \wmii. If
        prefixed with |xft:|, the Xft font renderer is used, and
        fonts may be antialiased.
      \item[grabmod ‹Modifier Keys›] The key which must be
        pressed to move and resize windows with the mouse
        without clicking hot spots.
      \item[incmode ‹Mode›] Controls how X11 increment hints are
        handled in managed mode. Possible values are:
        \begin{description}
          \item[ignore] Increment hints are ignored entirely.
            Clients are stretched to fill their full allocated
            space.
          \item[show] Gaps are shown around managed client
            windows when their increment hints prevent them from
            filling their entire allocated space.
          \item[squeeze] When increment hints cause gaps to show
            around clients, \wmii\ will try to adjust the sizes
            of the clients in the column to minimize lost space.
        \end{description}
      \item[view ‹Tag›] The currently visible view.
      \item[exec ‹Command›] Replaces this \wmii\ instance with
        ‹Command›. ‹Command› is split according to rc quoting
        rules, and no expansion occurs. If the command fails to
        execute, \wmii\ will respawn.
    \end{description}
  \item[props] The client's window class (the X11 |WM_CLASS|
    property) and title string, separated by colons. This file
    is not writable.
  \item[label] The client's window title. May be written to
    change the client's title.
  \item[tags] The client's tags. Tag names are separated by |+|
    signs. Tags beginning and ending with |/| are treated as
    regular expressions. If the written value begins with a |+|
    or a |-|, the tags are updated rather than overwritten. Tag
    names which directly follow a |-| sign are removed rather
    than added. Regular expression tags which directly follow a
    minus sign are treated as exclusion expressions. For
    example, the tag string |+/foo/-/food/| will match the tag
    |foobar|, but not the tag |foodstand|.
\end{description}

\subsection{Configuration}

We'll need to write our previously defined theme information to
\wmii:

\begin{code}
  «Configuration» ::=
  «Theme Definitions»

  xsetroot -solid $background
  wmiir write /ctl <<!
  border 2
  focuscolors $focuscolors
  normcolors $normcolors
  font $font
  grabmod Mod4
  !
\end{code}

\subsection{Key Bindings}

And we need a few more key bindings to select our views:

\begin{code}
  «Tag Selection Keys» ::=
  Mod4-Shift-t)
    # Prompt the user for a tag
    tags=$(wmiir ls /tag | sed ‘s,/,,; /sel/d’ | wimenu)
    # Write it to the filesystem.
    wmiir xwrite /ctl view $tag;;
  Mod4-[0-9])
    wmiir xwrite /ctl view ${2＃＃*-};;
\end{code}

\section{Tying it All Together}

\begin{code}
  #!/bin/sh
  «Click Menu Initialization»
  «Command Execution Initialization»

  «Configuration»

  «Bind Keys»
  «Event Loop»
\end{code}

\backmatter

\printindex

\end{document}
